\bigskip

Ceci peut permettre de penser que le susdit château de Nougaret pouvait bien se situer sur ce petit plateau délimité par l'actuel chemin de Saint André et la route de Villeneuve non loin du « Castellus » dont le nom de « CASTELLAS » rappellerait précisément l'existence de cette demeure féodale aujourd'hui disparue.

Pour mémoire, auraient été à cette époque « coseigneurs » de la baronnie d'ALBAN, en même temps que les « de Moret », les « de Vernhette » qui figurent sous ce titre de coseigneurs d'Alban dans une étude généalogique aveyronnaise. Notre documentation actuelle ne permet ni d'infirmer, ni de confirmer cette assertion... Il est vrai, près de trois siècles après, ce qui est bien loin, on trouve comme vicaire (curé) à Saint Michel de Saint Léon (aujourd'hui appelé St Michel de Léon), un Jean Vernhet.
